\cleardoublepage % memastikan bab baru dimulai di halaman ganjil (kanan)
\chapter{PENDAHULUAN}
\label{chap:pendahuluan}
% --- Latar Belakang ---
\section{Latar Belakang}
Algoritma dan Struktur Data (ASD) merupakan fondasi utama pendidikan ilmu komputer, namun berbagai studi menunjukkan bahwa mahasiswa konsisten mengalami kesulitan memahaminya karena sifat materinya yang abstrak, dinamis, dan multidimensional, yang menimbulkan beban kognitif tinggi serta memicu miskonsepsi sejak tahap awal pembelajaran \cite{Mtaho2024}. Kesulitan ini kerap bermula dari konsep dasar seperti \textit{pointer} dan struktur data linear, kemudian berkembang menjadi hambatan yang lebih serius pada topik lanjutan seperti \textit{tree}, \textit{graph}, dan analisis kompleksitas algoritma. Pembelajaran ASD di perguruan tinggi umumnya bertumpu pada modul tertulis, \textit{slide} presentasi, serta penjelasan sinkron dari dosen, yang efektif untuk penyampaian konsep terstruktur namun kurang mampu mengakomodasi kebutuhan belajar mandiri mahasiswa di luar jam kuliah, sehingga keterbatasan umpan balik yang cepat berkontribusi pada miskonsepsi yang sulit dikoreksi pada tahap selanjutnya \cite{Leinonen2023}.

Perkembangan \textit{Large Language Model} (LLM) membuka peluang baru dalam dunia pendidikan melalui kemampuannya menghasilkan penjelasan konsep, contoh implementasi kode, serta ilustrasi langkah-langkah algoritma secara dinamis dan kontekstual, sehingga berpotensi berperan sebagai teman belajar digital (\textit{AI learning companion}) yang mendukung pembelajaran mandiri yang lebih fleksibel \cite{Kasneci2023}. Meskipun demikian, penggunaan LLM secara langsung melalui \textit{chatbot} umum seperti ChatGPT atau Gemini belum cukup efektif sebagai sarana belajar yang terstruktur, karena interaksi percakapan bebas membuat penjelasan tidak selalu sesuai urutan materi kurikulum, dan LLM berpotensi menghasilkan informasi yang tidak akurat akibat fenomena \textit{hallucination} \cite{Raihan2024}. Tanpa mekanisme pengarahan yang terstruktur, keluaran LLM sulit dijamin kualitas dan relevansinya untuk keperluan pembelajaran akademik.

Untuk mengatasi keterbatasan tersebut, penelitian ini mengembangkan Teman Belajar ASD, sebuah aplikasi pembelajaran berbasis \textit{web} yang mengintegrasikan LLM dalam alur pembelajaran yang terkontrol. Pendekatan \textit{guided learning} menyediakan kerangka sistematis mulai dari pemahaman konsep, eksplorasi contoh, latihan soal, hingga ringkasan materi, sementara teknik \textit{prompt engineering} yang terstruktur membatasi keluaran LLM pada domain materi yang relevan, sehingga risiko \textit{hallucination} dan inkonsistensi konten dapat diminimalkan \cite{White2023PromptPatternCatalog}. Dengan demikian, diharapkan aplikasi ini dapat membantu mahasiswa memahami konsep ASD secara bertahap, mengurangi miskonsepsi, dan mendukung proses belajar mandiri yang lebih efektif.
% --- Rumusan Masalah ---
\section{Rumusan Masalah}
Secara umum, tugas akhir ini membahas pengembangan aplikasi pembelajaran berbasis \textit{web} yang memanfaatkan LLM sebagai teman belajar untuk menyajikan materi ASD secara terarah dan sesuai kurikulum. Berdasarkan uraian tersebut, rumusan masalah dalam tugas akhir ini adalah sebagai berikut:
\begin{enumerate}
    \item Bagaimana rancangan dan pengembangan aplikasi pembelajaran ASD berbasis \textit{web} dengan alur \textit{guided learning} yang terstruktur dan mudah diikuti oleh mahasiswa, sehingga dapat membantu mengurangi miskonsepsi pada konsep-konsep dasar ASD?

    \item Bagaimana memanfaatkan LLM dalam aplikasi sebagai teman belajar digital untuk menghasilkan penjelasan materi, contoh implementasi, dan latihan secara terkontrol agar tetap relevan, akurat, dan konsisten dengan kurikulum?
\end{enumerate}
% --- Tujuan ---
\section{Tujuan}
Berdasarkan rumusan masalah yang telah ditetapkan, tujuan penelitian ini adalah sebagai berikut:

\begin{enumerate}
    \item Mengembangkan aplikasi pembelajaran ASD berbasis \textit{web} dengan alur pembelajaran terstruktur menggunakan pendekatan \textit{guided learning} untuk membantu mahasiswa memahami konsep secara bertahap dan mengurangi miskonsepsi.

    \item Mengimplementasikan pemanfaatan LLM dalam aplikasi sebagai teman belajar digital yang mampu menghasilkan penjelasan materi, contoh implementasi, dan latihan secara relevan, akurat, serta konsisten dengan kurikulum.
\end{enumerate}
% --- Batasan Masalah ---
\section{Batasan Masalah}
Untuk memastikan ruang lingkup penelitian tetap terfokus dan sesuai dengan tujuan yang ingin dicapai, Tugas Akhir ini memiliki batasan-batasan sebagai berikut:

\begin{enumerate}
    \item Ruang lingkup materi yang digunakan dalam aplikasi dibatasi pada topik-topik dasar ASD seperti \textit{array}, \textit{linked list}, \textit{stack}, \textit{queue}, \textit{tree}, \textit{graph}, \textit{searching}, dan \textit{sorting} sesuai dengan kurikulum yang diberikan dosen.
    \item \textit{Large Language Model} (LLM) yang digunakan bersifat model siap pakai (\textit{pre-trained}) dan tidak mencakup proses pelatihan ulang (\textit{fine-tuning}) berskala besar. Penelitian ini hanya memanfaatkan teknik \textit{prompt engineering} dan \textit{retrieval-based guidance} untuk mengarahkan keluaran model.
    \item Interaksi pengguna utama dibatasi melalui mekanisme \textit{guided learning}. Sebagai fitur pelengkap, sistem menyediakan asisten percakapan kontekstual (\textit{chat asisten materi}) dan fasilitas evaluasi soal mandiri (\textit{soal sendiri}) yang memungkinkan mahasiswa mendapatkan penjelasan tambahan atau mengujikan soal dari sumber eksternal. Namun, seluruh interaksi tetap dibatasi pada konteks materi topik ASD yang sedang dipelajari dan tidak mencakup percakapan bebas di luar cakupan kurikulum.
    \item Sumber referensi yang digunakan LLM dibatasi pada dokumen yang disediakan dosen atau materi kuliah yang relevan. Sistem tidak membahas topik yang berada di luar cakupan kurikulum ASD.
    \item Evaluasi penelitian difokuskan pada kualitas konten dan kejelasan penyajian materi, serta aspek kegunaan aplikasi oleh mahasiswa, tanpa mencakup evaluasi performa teknis LLM secara mendalam.
\end{enumerate}

% --- Metodologi Pengerjaan TA ---
\section{Metodologi}
Metodologi yang digunakan dalam Tugas Akhir ini mengikuti kerangka \textit{Software Development Life Cycle} (SDLC) dengan tahapan berikut:

\begin{enumerate}
    \item \textbf{Perencanaan}

    Tahap ini mencakup identifikasi ruang lingkup penelitian, penentuan tujuan pengembangan aplikasi, serta penyusunan jadwal kerja. Aktivitas meliputi studi literatur terkait kesulitan mahasiswa dalam mempelajari ASD, pendekatan \textit{guided learning}, pemanfaatan LLM dalam pendidikan, serta penetapan tumpukan teknologi dan batasan implementasi yang sesuai dengan waktu pengerjaan.

    \item \textbf{Analisis}

    Pada tahap ini dilakukan analisis kebutuhan fungsional sistem, termasuk identifikasi aktor, alur pembelajaran yang diinginkan, dan kendala yang dialami mahasiswa saat belajar ASD secara mandiri. Analisis juga mencakup pemilihan model LLM yang digunakan (GPT-4o-mini melalui OpenAI API) serta perumusan strategi \textit{prompt engineering} untuk menjaga relevansi dan konsistensi keluaran model terhadap kurikulum.

    \item \textbf{Desain}

    Tahap desain berfokus pada perancangan arsitektur sistem berlapis (\textit{frontend}, \textit{backend}, dan \textit{database}), model basis data untuk menyimpan data pengguna dan progres belajar, serta alur \textit{guided learning} (materi $\rightarrow$ contoh $\rightarrow$ latihan $\rightarrow$ ringkasan). Selain itu, dirancang pula antarmuka halaman-halaman utama aplikasi dan \textit{template prompt} untuk setiap jenis konten yang dihasilkan LLM.

    \item \textbf{Implementasi}

    Implementasi dilakukan secara bertahap mengikuti rancangan yang telah disusun. \textit{Frontend} dibangun menggunakan Next.js dengan React, \textit{backend} menggunakan FastAPI berbasis Python, dan basis data menggunakan PostgreSQL. Integrasi LLM dilakukan melalui OpenAI API dengan GPT-4o-mini. Seluruh layanan di-\textit{deploy} di Google Cloud Run melalui \textit{pipeline} CI/CD berbasis GitHub Actions.

    \item \textbf{Pengujian}

    Tahap ini bertujuan untuk memastikan fungsionalitas sistem berjalan sesuai kebutuhan melalui pengujian berbasis skenario dengan pendekatan \textit{black box testing}. Evaluasi dilakukan pula terhadap kemudahan pengguna dalam mengikuti alur \textit{guided learning} melalui \textit{task-based usability testing}.
\end{enumerate}

% --- Sistematika Penulisan ---
\section{Sistematika Penulisan}
Laporan Tugas Akhir ini disusun secara sistematis agar memudahkan pembaca dalam memahami alur penelitian, perancangan, pengembangan, serta evaluasi aplikasi pembelajaran yang dibangun. Adapun sistematika penulisan dalam laporan ini adalah sebagai berikut.

\textbf{Bab I Pendahuluan} berisi latar belakang, rumusan masalah, tujuan penelitian, batasan masalah, metodologi yang digunakan, serta sistematika penulisan laporan secara keseluruhan.

\textbf{Bab II Studi Literatur} membahas landasan teoritis yang digunakan dalam penelitian ini, meliputi konsep pembelajaran Algoritma dan Struktur Data, karakteristik kesulitan mahasiswa dalam memahami ASD, prinsip kerja \textit{Large Language Model} (LLM), dan evaluasi. Bab ini juga menguraikan penelitian terdahulu yang relevan.

\textbf{Bab III Analisis} membahas kondisi pembelajaran Algoritma dan Struktur Data (ASD) saat ini, termasuk fungsi umum sistem pembelajaran yang digunakan mahasiswa serta berbagai kendala yang masih ditemui. Bab ini juga memuat analisis kebutuhan melalui identifikasi masalah pengguna, perumusan kebutuhan fungsional sistem, serta analisis alternatif solusi yang meliputi pemilihan model LLM dan mekanisme \textit{guided learning} yang paling sesuai untuk mendukung pengembangan aplikasi.

\textbf{Bab IV Perancangan} membahas perancangan sistem pembelajaran ASD berbasis \textit{web} menggunakan LLM, meliputi arsitektur sistem, alur kerja pembelajaran, perancangan basis data, serta visualisasi antarmuka aplikasi yang akan diimplementasikan.

\textbf{Bab V Implementasi} berisi penjelasan tentang proses implementasi solusi yang diusulkan, termasuk langkah-langkah yang diambil, teknologi yang digunakan, dan hasil yang diperoleh.

\textbf{Bab VI Evaluasi} berisi analisis dan evaluasi terhadap solusi yang diimplementasikan, termasuk metode pengujian, hal-hal yang disiapkan sebelum pengujian, hasil pengujian, dan penilaian kinerja.

\textbf{Bab VII Kesimpulan dan Saran} berisi kesimpulan dari hasil pelaksanaan tugas akhir dan saran untuk pengembangan lebih lanjut.
